\input{preamble.tex}
\title{Coset Helly properties in an HNN extension}
\author{Human Compression project}
\date{Private draft, 11 October 2026}
\begin{document}
\maketitle

Let $A$ be a group and let $\mathcal P$ be a family of its subgroups,
closed under finite intersections. We say that $\mathcal P$ has the
\emph{coset Helly property} if every finite family of pairwise intersecting
left cosets of members of $\mathcal P$ has nonempty total intersection.
We include $A$ in $\mathcal P$. All standard parabolics of a finite-rank
flag-transitive coset geometry have this property: replace a coset of
an intersection by the corresponding family of maximal-parabolic cosets,
then use flag transitivity. Multiple cosets of the same maximal parabolic,
if they intersect, coincide, so repeated types cause no problem.

Fix $C,D\in\mathcal P$ and an isomorphism $\phi:C\longrightarrow D$.
Write
\[
 G=\langle A,t\mid t^{-1}ct=\phi(c)\ (c\in C)\rangle.
\]
A subgroup $P\in\mathcal P$ is called compatible if
$\phi(P\cap C)=P\cap D$. Put $\widehat P=\langle P,t\rangle$.
Britton's lemma embeds the HNN extension of $P$ along these intersections
in $G$ and gives $\widehat P\cap A=P$. These are retained classical
normal-form facts. Unlike an arbitrary family of special subgroups, the
family below keeps its local coset Helly hypothesis.

\begin{theorem}\label{thm:helly}
If $\mathcal P$ has the coset Helly property, then the family consisting
of all $P\in\mathcal P$ and all $\widehat P$ for compatible
$P\in\mathcal P$ has the coset Helly property in $G$.
For compatible $P,Q$,
\[
 \widehat P\cap\widehat Q=\widehat{P\cap Q},\qquad
 R\cap\widehat P=R\cap P\quad(R\in\mathcal P).
\]
\end{theorem}
\begin{proof}
Use the Bass--Serre tree $T$ of this HNN extension. Its vertices are
$gA$ and its oriented edges are $gC$, with endpoints $gA$ and $gtA$.
The latter endpoint is independent of the representative because
$ct=t\phi(c)$. At a vertex $xA$, outgoing edge directions are
parametrized by left cosets of $C$ in $A$, and incoming directions by
left cosets of $D$ in $A$.

For compatible $P$, the embedded Bass--Serre tree of $\widehat P$ has
vertices $hA$ and edges $hC$, $h\in\widehat P$. Its injectivity follows
from the same reduced word theorem, since the associated intersections
are exactly $P\cap C$ and $P\cap D$. Denote this connected subtree by
$T_P$. For a base subgroup $R\leq A$ instead use the singleton subtree
$\{A\}$. Attach to a coset $g\widehat P$ the subtree $gT_P$ and to a
coset $gR$ the vertex $gA$. These assignments are independent of the
chosen coset representative.

If a vertex $v=xA$ belongs to such a support subtree, the corresponding
coset meets $xA$. After multiplication by $x^{-1}$ this intersection is
a left coset in $A$ of $P$, or of $R$ in the base-subgroup case.
For example, choose $y\in g\widehat P\cap xA$; the intersection is
$y(\widehat P\cap A)=yP$. An incident edge belongs to $gT_P$ precisely
when this local $P$ coset meets the local $C$ or $D$ coset describing
the direction of that edge. This follows directly from the vertex and
edge descriptions of the embedded tree.

We need one transport observation. Suppose two global cosets meet, and
a vertex $v$ lies in both support subtrees. Choose a common element $z$
of the cosets. Both support subtrees contain $zA$, so they contain the
whole finite path from $zA$ to $v$. Start with the common representative
$z$ at its first vertex. At the next edge, the two local subgroup cosets
intersect each other and each meets the local associated-subgroup coset
of that edge. Local coset Helly supplies a common element of all three.
For an outgoing edge multiply this representative on the right by $t$;
for an incoming edge multiply it by $t^{-1}$. The result lies at the
next vertex and still belongs to both global cosets, because both
subgroups contain $t$. If either subgroup was a base subgroup, its
support was a singleton and the path has length zero, so no multiplication
is needed. Repeating the step gives a common representative in $xA$.

Now take finitely many pairwise intersecting global cosets. Their support
subtrees intersect pairwise. The elementary Helly property for subtrees
gives a common vertex $v=xA$. By the transport observation every pair
of local fibers at $v$ intersects. These fibers are cosets of members
of $\mathcal P$, so local coset Helly supplies an element in all of
them. Multiplying by $x$ proves the first assertion.

For the intersection formula, first observe that $P\cap Q$ is compatible:
an isomorphism commutes with intersections. Thus
$K=\widehat{P\cap Q}$ is contained in both groups. Take
$g\in\widehat P\cap\widehat Q$ and follow the finite path from $A$
to $gA$. Both embedded subtrees contain it. We show inductively that
every vertex on this path is $kA$ for some $k\in K$.
This holds initially with $k=1$. Translate a current vertex $kA$ and
its next edge by $k^{-1}$. If that edge is outgoing, its local coset
$aC$ meets both $P$ and $Q$; if incoming, its local coset $aD$ does.
Since $P\cap Q$ contains $1$, the three local cosets meet pairwise.
The local Helly property supplies $p\in P\cap Q$ in the relevant
coset. The next vertex is therefore $kp tA$ or $kp t^{-1}A$,
respectively, and its representative is again in $K$.
At the end $gA=kA$, and hence
$k^{-1}g\in A\cap\widehat P\cap\widehat Q=P\cap Q\leq K$.
This proves the first equality. The second follows immediately from
$\widehat P\cap A=P$.
\end{proof}

\begin{corollary}\label{cor:geometry}
Let $(A,(A_i)_{i\in I})$ be a flag-transitive coset geometry of finite
rank. Let $C,D$ be standard parabolics, and let $\phi:C\to D$ be an
isomorphism. Choose a partition $\mathcal S$ of $I$ such that every
$A_s=\bigcap_{i\in s}A_i$, $s\in\mathcal S$, is compatible, and choose
$F\subseteq I$. In the HNN extension take maximal parabolics
\[
 G_s=\langle A_s,t\rangle\quad(s\in\mathcal S),\qquad
 G_*=A_F.
\]
This is a flag-transitive coset geometry. For a set $J$ of its types,
write $E=\bigcup_{s\in J\cap\mathcal S}s$. Then
\[
 G_J=
 \begin{cases}
   \langle A_E,t\rangle,&*\notin J,\\
   A_{E\cup F},&*\in J.
 \end{cases}
\]
Its chamber stabilizer is $A_I$.
\end{corollary}
\begin{proof}
Take $\mathcal P$ to be all standard parabolics of the input, including
$A$. These have the local coset Helly property explained above. Theorem
\ref{thm:helly} applies to all the displayed maximal parabolics. A flag
is a finite family of pairwise intersecting such cosets, so it has a
common representative. It is therefore a translate of a standard flag
and extends to a standard chamber. The intersection formulas follow
by induction from the theorem, with the usual empty-intersection convention.
For the complete type set $E=I$, so the Borel is $A_I$.
\end{proof}

\begin{corollary}\label{cor:rc}
Suppose the input in Corollary~\ref{cor:geometry} is residually connected.
Let $C=A_J$, $D=A_K$, and suppose a specified bijection
$f:I\setminus J\longrightarrow I\setminus K$ satisfies
\[
 \phi(A_{J\cup X})=A_{K\cup f(X)}\qquad(X\subseteq I\setminus J).
\]
Take $\mathcal S$ to be the connected classes of the arrows $i\to f(i)$
on $I$, and put $F=J\setminus K$. Then the resulting HNN geometry is
residually connected. Its $*$ panels are infinite. Its panel of type
$s\in\mathcal S$ has size
\[
 [A_{I\setminus(s\setminus F)}:A_I].
\]
In particular it is firm when the input is firm.
\end{corollary}
\begin{proof}
Write $H=A_I$ and $R_i=A_{I\setminus\{i\}}$.
The displayed compatibility gives $\phi(H)=H$ and
$\phi(R_i)=R_{f(i)}$ for $i\notin J$.
Each class $s$ satisfies $f(s\setminus J)=s\setminus K$.
Consequently $\phi(A_s\cap C)=A_s\cap D$, so the previous corollary
applies. Its minimal parabolics are
\[
 M_s=A_{I\setminus(s\setminus F)},\qquad M_*=\langle H,t\rangle.
\]

Residual connectivity of the input gives the elementary generation rule
\[
 A_{I\setminus E}=\langle H,R_i\ (i\in E)\rangle
 \quad(E\subseteq I).
\]
For $E=\varnothing$ this just retains the Borel $H$; for nonempty $E$
it follows by induction from generation by maximal parabolics in each
connected residue of rank at least two. This is the Borel-preserving
version of the usual residual generation criterion.

Let $W$ be a subset of the new type set $\mathcal S\cup\{*\}$, and
put $E=\bigcup_{s\in W\cap\mathcal S}s$. If $*\notin W$, then
\[
 \langle H,M_s\ (s\in W)\rangle
 =A_{I\setminus(E\setminus F)}=G_{(\mathcal S\cup\{*\})\setminus W}.
\]
If $*\in W$, this generated group also contains $t$.
Every $i\in E\cap F$ is in the range but not the domain of $f$.
It therefore has a predecessor $j\in E\setminus F$ with $f(j)=i$.
Conjugation by $t$ brings $R_i$ into the group from $R_j$.
Thus in this case the generated group is
$\langle H,t,R_i\ (i\in E)\rangle=\langle A_{I\setminus E},t\rangle$,
again exactly $G_{(\mathcal S\cup\{*\})\setminus W}$.

We have proved the same generation rule for every new standard parabolic.
If at least two types remain, each of its minimal generators belongs to
one of the parabolics obtained by deleting a different remaining type.
Those parabolics therefore generate it. The residual generation criterion
proves residual connectivity.

Since $\phi(H)=H$, the subgroup $\langle H,t\rangle$ is an extension
of $H$ by the infinite cyclic group generated by the stable letter.
Its index over $H$ is infinite. The other panel indices are the stated
indices of $M_s$ over $H$. Every class has a point outside $F$: a point
in $F$ has a predecessor outside $F$. If the input is firm, the group
$M_s$ therefore contains some $R_i$ strictly larger than $H$.
All these panels have at least two elements, as required.
\end{proof}

The partition and compatibility conditions in Corollary~\ref{cor:geometry}
are explicit. No isomorphism of unlabelled subgroup lattices is silently
promoted to an action on named types. If a specified partial bijection
$I\setminus J\to I\setminus K$ describes an isomorphism
$A_J\to A_K$ on the indexed parabolics, its connected type classes
give the partition used in Piedade--Tranchida, arXiv:2505.24662v1,
Section 3.3. Verifying this precise condition recovers their forward
flag-transitivity and intersection assertions. Corollary~\ref{cor:rc}
supplies residual connectivity, firmness and exact panel sizes under
the explicit indexed compatibility.

The proof combines classical Bass--Serre normal forms with the elementary
Helly property for subtrees and local coset Helly. The transfer argument
is written here to repair the use of the unrestricted special-subgroup
intersection lemma, which fails in the separate obstruction note.
Historical priority for this transfer formulation has not been assessed.
This manuscript is a private draft awaiting distinct review; no numerical
or complete-original credit is assigned.
\end{document}
